% ============================================================
% CONDITIONAL CANON LEMMA v15B.0: single-modulus resolved core envelope
% ============================================================
\subsection{Conditional Single-Modulus Core-Envelope Lemma}
\label{subsec:canon_single_modulus_core_envelope}

\paragraph{Status.}
\textbf{[CANON-COMPATIBLE CONDITIONAL LOCAL CORE LEMMA].}  If the resolved internal SST core
envelope is represented, in the transverse core plane, by one complex local
order parameter with a single isotropic stiffness modulus.  This lemma applies
only to the internal core-envelope layer; it does not identify the internal
core sound speed with the transverse causal speed of the torsion/shear layer.

Let
\begin{equation}
    \Psi(\rho,\theta)=F(\rho)e^{in\theta},
    \qquad n\in\mathbb Z,
\end{equation}
be the local resolved core envelope.  The canonical minimal transverse energy
is
\begin{equation}
    E_\perp[\Psi]
    =
    \kappa\int_{\mathbb R^2}
    \left[
        |\nabla\Psi|^2
        +\frac{1}{2\xi^2}(1-|\Psi|^2)^2
    \right]d^2x .
    \label{eq:canon_single_modulus_core_energy}
\end{equation}
Under this local-core assumption, the single-modulus statement is the equality of the radial-amplitude and
azimuthal-phase stiffnesses inside the Euclidean transverse norm
\(|\nabla\Psi|^2\).  It is not a statement about the external radiation layer.

In polar coordinates,
\begin{equation}
    |\nabla\Psi|^2
    =
    \left(\frac{dF}{d\rho}\right)^2
    +\frac{n^2F^2}{\rho^2} .
\end{equation}
With \(\rho=\xi r\), Eq.~\eqref{eq:canon_single_modulus_core_energy} reduces,
up to the common factor \(2\pi\kappa\), to
\begin{equation}
    E_\perp
    \propto
    \int_0^\infty
    \left[
        F'^2r
        +n^2\frac{F^2}{r}
        +\frac12(F^2-1)^2r
    \right]dr .
    \label{eq:canon_dimensionless_gp_core_energy}
\end{equation}
Comparing Eq.~\eqref{eq:canon_dimensionless_gp_core_energy} with the generic
radial density
\begin{equation}
    \mathcal L_r
    =
    A F'^2 r
    +B n^2\frac{F^2}{r}
    +\frac{C}{2}(F^2-1)^2r
\end{equation}
gives the internal coefficient identity
\begin{equation}
    \boxed{A=B=C.}
    \label{eq:canon_internal_abc_identity}
\end{equation}
The Euler--Lagrange equation for the generic form is
\begin{equation}
    F''+\frac{F'}{r}
    -\frac{B}{A}n^2\frac{F}{r^2}
    +\frac{C}{A}F(1-F^2)=0 .
\end{equation}
Therefore, in the single-quantum sector \(n=1\), the canonized core-envelope
lemma gives
\begin{equation}
    \boxed{
    F''+\frac{F'}{r}-\frac{F}{r^2}+F(1-F^2)=0 .
    }
    \label{eq:canon_unit_gp_nlse_core_ode}
\end{equation}
Consequently the depletion-energy term in the dimensionless radial functional
must carry coefficient \(1/2\).  A coefficient \(1/4\) would imply
\(C/A=1/2\) and would not vary to Eq.~\eqref{eq:canon_unit_gp_nlse_core_ode}.

\paragraph{Canon boundary.}
This lemma records the local one-modulus internal core-envelope reduction and
its consequence \(A=B=C\) as a conditional theorem.  The remaining gate is an
independent derivation or explicit canon acceptance of the single-modulus
Madelung core envelope itself.  It does not canonize an exact identification of
\(\alpha_{\rm ring}\) with the legacy value \(1.61\), with \(\varphi\), or with
the electromagnetic fine-structure constant.  Nor does it promote the
Lorentz-compatible Swirl-Clock to a theorem from the NLSE core alone; the
separate core--torsion impedance gate remains required for the external causal
layer.
